% Figure from MATH 590 notes, section 15. Compile with the notes preamble.
\begin{tikzpicture}[x=6.2cm, y=1cm]
  \draw[black!40, thin] (-0.36,0) -- (1.12,0);
  % U contains (-eps, eps) cap X
  \fill[figR!15] (-0.28,-0.18) rectangle (0.28,0.18);
  \draw[figR, thick, dashed] (-0.28,-0.18) rectangle (0.28,0.18);
  \node[font=\scriptsize, figR] at (0,0.5) {$U \supseteq (-\varepsilon,\varepsilon)\cap X$};
  \node[font=\scriptsize, figR] at (-0.28,-0.45) {$-\varepsilon$};
  \node[font=\scriptsize, figR] at (0.28,-0.45) {$\varepsilon$};
  % the tail 1/n, n >= 4, all inside U
  \foreach \n in {4,...,40} \fill[figB] ({1/\n},0) circle (0.9pt);
  \fill[figB] (0,0) circle (1.4pt);
  \node[font=\scriptsize, below=3pt] at (0,0) {$0$};
  % finitely many points left over, each in its own cover element
  \foreach \n/\val in {3/{\tfrac13}, 2/{\tfrac12}, 1/{1}} {
    \pgfmathsetmacro{\p}{1/\n}
    \draw[black!55, thick] (\p-0.028,0.3) -- (\p+0.028,0.3);
    \draw[black!55, fill=white, thick] (\p-0.028,0.3) circle (1.2pt);
    \draw[black!55, fill=white, thick] (\p+0.028,0.3) circle (1.2pt);
    \fill[figB] (\p,0) circle (1.4pt);
    \node[font=\scriptsize, below=3pt] at (\p,0) {$\val$};
  }
  \node[font=\scriptsize, black!70] at (1/3,0.62) {$U_{y_3}$};
  \node[font=\scriptsize, black!70] at (1/2,0.62) {$U_{y_2}$};
  \node[font=\scriptsize, black!70] at (1,0.62) {$U_{y_1}$};
  \node[font=\scriptsize, figB] at (0.72,-0.95) {leftover points $y_1,y_2,y_3$};
\end{tikzpicture}
